\documentclass[11pt]{article}

\usepackage{amsmath, amssymb}
\usepackage{geometry}
\usepackage{hyperref}
\usepackage{booktabs}
\usepackage{microtype}
\usepackage{parskip}

\geometry{margin=1.1in}

\hypersetup{
    colorlinks=true,
    linkcolor=blue,
    urlcolor=blue
}

\title{\textbf{CTMC Models for Customer Journey Prediction}\\[0.4em]
\large UCLA Stats 148 Capstone --- Technical Reference}
\author{M148 Project Team}
\date{May 2026}

\begin{document}
\maketitle
\tableofcontents
\newpage

\section{Problem Setup}

Each customer journey is a timestamped sequence of events. The project label is
binary:
\[
    y =
    \begin{cases}
    1 & \text{journey reaches \texttt{order\_shipped}, \texttt{ed\_id}=28,} \\
    0 & \text{journey fails by taking no customer action for 60 days.}
    \end{cases}
\]

This is not exactly the same target as the original CTMC experiments. The first
CTMC models estimated
\[
    P(T_{\text{success}} \leq 60\text{ days} \mid X(0)=s),
\]
the probability of hitting the success state within a fixed horizon. The
classification target is better described as
\[
    P(T_{\text{success}} < T_{\text{inactivity-failure}} \mid X(0)=s),
\]
where the failure time occurs when no customer action happens for 60 days.

\section{State-Space Choice}

The file \texttt{data/Event Definitions.csv} mixes customer actions with
company or system events. For the inactivity definition, company/system events
should not reset the 60-day clock. Therefore the active CTMC state space uses
customer-initiated actions only:
\[
\{2,3,4,5,6,7,8,9,10,11,18,19,22,23,25\}.
\]
The success event \texttt{order\_shipped} (\texttt{ed\_id}=28) is retained as a
terminal outcome state, not as a customer action. The timeout-aware models also
add an artificial absorbing failure state \texttt{-1}.

This design intentionally excludes states such as approval, decline, pending,
catalog mail, downpayment received/cleared, fraud review, account activation,
and order fulfillment events from the behavioral CTMC state space.

\section{Retained Baseline 1: Global CTMC}

\subsection{Fitting}

For consecutive customer-action transitions $(i,j,\Delta t)$, define
\[
    N_{ij} = \#\{(i,j,\cdot)\}, \qquad
    T_i = \sum_{(i,j,\Delta t)} \Delta t.
\]
The maximum-likelihood CTMC generator is
\[
    \hat q_{ij} = \frac{N_{ij}}{T_i}, \quad i \neq j,
    \qquad
    \hat q_{ii} = -\sum_{j \neq i}\hat q_{ij}.
\]

\subsection{Prediction Target}

The global baseline still predicts the old target:
\[
    \hat p_{\text{global}}(s)
    =
    \left[\exp(Q^* \tau)\right]_{s,28},
    \qquad
    \tau = 60 \cdot 86400.
\]
Here $Q^*$ makes state 28 absorbing. This is useful as a baseline but it does
not explicitly model inactivity failure.

\section{Retained Baseline 2: Clustered CTMC}

The clustered CTMC groups journeys by prefix behavior and fits one global CTMC
per cluster:
\[
    Q^{(k)} = \texttt{GlobalCTMC.fit}(\mathcal{D}_k).
\]

The implementation supports two clustering modes:

\begin{center}
\begin{tabular}{lll}
\toprule
Mode & Switch & Notes \\
\midrule
KMeans & default & Original baseline; clusters action proportions and state features. \\
Spectral & \texttt{use\_spectral\_clustering=True} & Nearest-neighbor spectral fit plus KNN assignment for open journeys. \\
\bottomrule
\end{tabular}
\end{center}

The clustered CTMC also predicts the old success-hit target unless paired with
one of the timeout-aware models below.

\section{New Model 1: Timeout-Absorbing CTMC}

\subsection{Construction}

The timeout transition table modifies the training data:

\begin{itemize}
    \item Customer-action transitions are kept.
    \item Transitions into \texttt{order\_shipped}=28 are kept as success.
    \item Each incomplete journey receives one artificial terminal transition
          from its last customer-action state to failure state $f=-1$ after
          60 days.
\end{itemize}

A CTMC generator is then fit on this augmented state space. Both success and
failure are absorbing.

\subsection{Prediction}

For transient states $\mathcal{T} = \mathcal{S}\setminus\{28,-1\}$, the
probability of eventually absorbing in success before failure solves
\[
    Q_{\mathcal{T}\mathcal{T}}u
    =
    -Q_{\mathcal{T},28}.
\]
The model returns $u_s$ for the current state $s$. This directly targets
\[
    P(T_{\text{success}} < T_{\text{failure}} \mid X(0)=s).
\]

\section{New Model 2: Empirical Semi-Markov Timeout Model}

The semi-Markov timeout model removes the exponential waiting-time assumption.
For each state $i$, it stores empirical outgoing samples
\[
    (j_m, \Delta t_m).
\]
The value function is updated by
\[
    V(i) =
    \frac{1}{M_i}\sum_{m=1}^{M_i}
    \begin{cases}
        1, & j_m=28 \text{ and } \Delta t_m \leq 60\text{ days}, \\
        0, & j_m=-1 \text{ or } \Delta t_m > 60\text{ days}, \\
        V(j_m), & \text{otherwise.}
    \end{cases}
\]
Iteration continues to convergence. This estimates success before timeout
using the empirical waiting-time distribution rather than CTMC exponential
holding times.

\section{Active Model Set}

The active CTMC comparison now contains exactly two old CTMC methods and two
timeout-aware methods:

\begin{center}
\begin{tabular}{lll}
\toprule
Model & Target & Status \\
\midrule
\texttt{GlobalCTMC} & hit success within 60 days & retained old baseline \\
\texttt{ClusteredCTMC} & hit success within 60 days & retained old baseline \\
\texttt{TimeoutAbsorbingCTMC} & success before inactivity failure & new primary CTMC target \\
\texttt{SemiMarkovTimeoutModel} & success before inactivity failure & new non-exponential target \\
\bottomrule
\end{tabular}
\end{center}

Older experimental CTMC variants such as neural-rate CTMC, XGBoost-rate CTMC,
time-stratified CTMC, Laplace-smoothed CTMC, and isotonic CTMC calibration are
not part of the active CTMC comparison.

\section{Commands}

Run the retained CTMC comparison:
\[
\texttt{python -m src.models.ctmc\_pipeline all --no-cache}
\]

Run only spectral clustered CTMC:
\[
\texttt{python -m src.models.ctmc\_pipeline clustered --use-spectral-clustering --no-cache}
\]

\end{document}
